\section{Threat Model and Task Definition}
\label{sec:task}

\paragraph{Setting.}
A user issues a request $q$ to an LLM application. To answer it, the application places an external context $c$ (a
document, a table, code, or one tool output of an agent) into the model's input. The user and the application
developer are trusted. The author of $c$ is not: an attacker may have written part of $c$ to make the model do
something the user did not ask for. The attacker controls only the text of $c$, cannot see or change $q$, and does not
adapt to the detector; our attack strings are static. A detector $f$ sits in front of the model and maps an input to
a score $f(q, c) \in [0,1]$; above a threshold $\tau$ the context is flagged and withheld. Detectors that read only the
context compute $f(c)$.

\paragraph{Two definitions of an injection.}
Under the \emph{strict} view, every instruction in $c$ is untrusted, whatever $q$ asks for
\citep{chen2024struq}. The label then depends on $c$ alone: $y(q,c) = \mathbf{1}[c \text{ contains an
instruction}]$. Under the \emph{intent-based} view that we adopt, an instruction $i$ in $c$ is an injection
\emph{relative to the request}: it is one if following it while answering $q$ would give the user something other
than what they asked for, and it is not one if following it serves $q$. This makes BIPIA's ``deviation from user
expectations'' \citep{yi2023bipia} request-specific, and matches the aligned/misaligned distinction of
\citet{jia2026alignsentinel} and the conflict-based framing of the instruction hierarchy
\citep{wallace2024hierarchy}. The label is $y(q, c) = \mathbf{1}[\exists\, i \in c:\ i \text{ subverts } q]$.

\paragraph{Why the distinction matters for detectors.}
The two views agree on malicious payloads, which subvert every plausible request, and on contexts with no
instruction. They disagree on instructions that the user requested or that serve the request (an email that says
how to format the reply the user asked for, a README step the user wants carried out). A content-only detector
computes the same score for $(q_1, c)$ and $(q_2, c)$, so on any set where the label of $c$ changes with $q$ it is
at chance. Under the strict view this is no loss, and flagging requested instructions is correct. Under the
intent-based view it is a false alarm, and with agents that act on their tool outputs it blocks requested work.

\paragraph{Scope of our claims.}
We state the intent-based definition as a design choice, not as the only correct one. Which of our results depend on
it is fixed in advance: the pair-based hypotheses (H2, H3 on pairs, and H4 on false alarms for requested
instructions) are meaningful only under the intent-based view, whereas the artifact analysis, the \vthree{} gates,
and transfer to agent benchmarks (H1, and H3 on content-only sets) are meaningful under both. To keep the contested
zone small, our pairs exclude every instruction that would be unsafe to follow even when the request asks for it:
sending or revealing data to a party the user did not name, payments, deleting or overwriting user data, changing
credentials or permissions, contradicting the document, and unrequested promotional content (\S\ref{sec:pairs}).
Under either view those are injections.

\paragraph{Task format.}
Every model we train sees \inp{$q$}{$c$}; the cross-encoder of \S\ref{sec:ict} receives $q$ and $c$ as a segment
pair. Public content-only detectors see $c$ alone. On agent benchmarks, $q$ is the task's user prompt and $c$ is a
single tool output.

\paragraph{Metrics.}
Headline metrics are the true-positive rate at 1\% false-positive rate (TPR@1\%FPR), AUROC, and the false-positive
rate at the detector's default threshold on clean inputs. Thresholds for TPR@FPR are primarily \emph{calibrated}:
fixed on a held-out clean set and applied unchanged, as a deployment would. The \emph{oracle} protocol, which picks the
threshold on the evaluated set's own clean rows as \citet{liu2026passing} does, is reported next to it and never
decides a hypothesis. TPR@1\%FPR is reported only where the calibration set has at least 500 clean rows, and
TPR@5\%FPR only where it has at least 100; otherwise we report n/a. On pairs, the metric is pair accuracy: the
fraction of pairs in which the subverting member scores higher than the aligned one, with ties counted as one half.
Confidence intervals are 95\% cluster bootstraps ($B{=}1000$) that resample source documents (or user tasks on agent
benchmarks), and differences between models use a paired cluster bootstrap.
